\section{全书总架构：从零理解 AI 编译器、推理引擎与 C++20 后端}

\subsection{先把“编译器”讲成人话}

如果完全没有学过编译原理，可以先把编译器理解成一个严谨的翻译和审计系统。它不只是把一种文本换成另一种文本，而是要回答三个问题：

\begin{enumerate}
  \item 这份输入到底表达了什么语义。
  \item 这个语义能不能在目标机器上安全、正确地执行。
  \item 在不改变语义的前提下，怎样让机器少做无用工作。
\end{enumerate}

普通 C++ 编译器的输入是源代码。它会先把字符切成 token，再按语法组成 AST，然后检查类型、名字、生命周期、调用规则和错误边界，再把程序降成 IR，做优化，最后生成目标机器代码。AI 编译器的输入不一定是 C++ 源码，而是模型配置、tokenizer、权重文件、量化元数据和一张推理图。它要做的仍然是同一件事：先理解语义，再建立可检查的中间表示，最后把它降到 CPU/GPU 后端能执行的 kernel 计划。

\begin{keyidea}
编译原理不是一套抽象名词，而是一种工程纪律：任何优化都必须先知道输入语义、表示边界、错误条件和目标机器成本。AI 编译器只是把这套纪律用在模型、张量、量化和推理运行时上。
\end{keyidea}

本册不会假设读者已经懂词法分析、语义分析、IR、优化、lowering、runtime 这些词。每次使用它们时，都要落到本地推理引擎里的具体对象。例如 manifest 就像模型文件的符号表；shape verifier 就像类型检查器；graph IR 就像算子级中间表示；memory planner 就像给大张量做寄存器/栈槽分配；backend lowering 就像把 IR 变成 AVX、线程池、GPU stream 和 kernel launch。

\subsection{第三册怎样承接第一册和第二册}

第三册不是另开一条 AI 支线，而是把前两册的能力集中用在 AI infra 上。第一册解决“一个 C++ 程序怎样落到机器事实”：源码怎样编译、汇编和链接，进程怎样获得虚拟地址空间，寄存器、栈、对象布局、指针和函数调用怎样受 ABI 约束，怎样用反汇编、调试器和 benchmark 证明自己运行的是哪份二进制。没有这部分基础，后端优化里的 \code{std::span} 张量视图、汇编形态、调用边界、aligned buffer、C ABI kernel wrapper 和可信测量都会变成空话。

第二册解决“同一段机器工作进入现代计算系统以后怎样变慢或变快”：取指前端、乱序执行、分支预测、cache、TLB、page fault、SIMD、线程、锁、原子、false sharing、NUMA、任务运行时、I/O、背压和分布式边界怎样改变成本。没有这部分基础，第三册里讨论 GEMM/GEMV、packing、KV cache、prefill/decode、线程池、NUMA 放置、GPU staging 和 profiler 时，只能停在“经验上这样更快”。

\begin{center}
\begin{tabular}{p{0.18\linewidth}p{0.34\linewidth}p{0.34\linewidth}}
\toprule
卷次 & 已建立能力 & 第三册中的用法 \\
\midrule
第一册 & C++ 到二进制、汇编、ABI、对象布局、可信测量 & 读 kernel 汇编、设计 C++20 RAII/ABI 边界、验证 benchmark 身份 \\
第二册 & 现代 CPU、cache/TLB、SIMD、多线程、NUMA、运行时、背压 & 解释算子瓶颈、设计 packing/threading、分析 KV cache 和 prefill/decode \\
第三册 & 模型、张量、AI 编译器、推理 runtime、CPU/GPU 后端 & 把前两册能力组合成可验证的本地大模型推理引擎 \\
\bottomrule
\end{tabular}
\end{center}

因此，第三册每个重要机制都要向前接两条线：语义和 ABI 线来自第一册，硬件和系统成本线来自第二册。比如一个 int8 GEMV kernel，第一册帮助我们看懂函数边界、指针、汇编和测量；第二册帮助我们解释 SIMD 宽度、cache line、TLB、线程划分和 NUMA；第三册再把它放进模型权重、量化 scale、typed graph、runtime segment 和 oracle/profiler 中。

\subsection{AI infra 的五层：从文件到 token 流}

一台真正可解释的本地大模型推理引擎，可以分成五层。每一层都有自己的输入、输出、错误边界和优化空间。

\begin{systemdiagram}{第三册的五层 AI infra 架构}
\begin{pdfdiagram}[x=1cm,y=1cm]
  \node[book accent node,text width=2.35cm] (file) at (0,2.8) {模型资产层\\config/tokenizer\\weights};
  \node[book teal node,text width=2.35cm] (front) at (2.9,2.8) {编译前端\\manifest\\verifier};
  \node[book purple node,text width=2.35cm] (ir) at (5.8,2.8) {IR 与优化层\\graph/tensor\\fusion};
  \node[book amber node,text width=2.35cm] (backend) at (8.7,2.8) {后端层\\CPU/GPU\\kernel};
  \node[book navy node,text width=2.35cm] (rt) at (11.6,2.8) {运行时层\\session/KV\\sampler};

  \node[book label,text width=2.35cm] at (0,1.25) {字节是否完整\\格式是否可信};
  \node[book label,text width=2.35cm] at (2.9,1.25) {shape/dtype\\量化合同};
  \node[book label,text width=2.35cm] at (5.8,1.25) {语义不变\\减少内存流};
  \node[book label,text width=2.35cm] at (8.7,1.25) {SIMD/线程\\显存/stream};
  \node[book label,text width=2.35cm] at (11.6,1.25) {prefill/decode\\证据闭环};

  \draw[book strong arrow] (file) -- (front);
  \draw[book strong arrow] (front) -- (ir);
  \draw[book strong arrow] (ir) -- (backend);
  \draw[book strong arrow] (backend) -- (rt);
\end{pdfdiagram}
\diagramnote{本图要证明的命题是：AI infra 不是一个大函数，而是模型资产、编译前端、IR 优化、后端和运行时五层合同。}
\end{systemdiagram}

\topic{模型资产层。} 这里处理的是文件事实：模型配置写了多少层，tokenizer 有多少词，权重张量在哪个文件偏移，量化 scale 怎样存储。错误通常是文件缺失、shape 不一致、tokenizer 和 embedding 不匹配、量化 metadata 不完整。

\topic{编译前端层。} 这里把文件事实变成编译器事实。manifest 记录张量和模型合同，schema resolution 把不同模型命名映射成统一角色，shape/dtype/quant verifier 拒绝不可信输入。前端做得越扎实，后端越不需要在热路径里猜。

\topic{IR 与优化层。} 这里把推理图变成 graph IR、tensor IR 和 loop/kernel IR。Graph IR 负责算子依赖和状态边；Tensor IR 负责 shape、stride、layout 和 dtype；Loop IR 负责 tile、vector、thread 和 memory access。融合、消除冗余 reshape、layout 选择和内存生命周期分析都发生在这里。

\topic{后端层。} 这里把计划落到机器。CPU 后端要处理 cache line、SIMD、packing、线程划分、NUMA、TLB 和分支；GPU 后端要处理显存布局、coalesced load、shared memory、kernel launch、stream 和同步。后端优化是本册后半部分的重头戏，不能停留在“调用一个库”。

\topic{运行时层。} 编译器生成计划后，runtime 还要维护 session、token buffer、KV cache、activation arena、sampler、profiler 和 oracle。prefill 和 decode 是不同阶段，指标也不同：prefill 关心首 token 和 prompt 吞吐，decode 关心 tokens/s、单 token latency 和尾延迟。

\subsection{整本书按六个大章组织}

本册后续按六个大章推进。每一节都回答一个具体问题，而不是堆主题名。

\topic{第一章：总架构与业务闭环。} 先把 AI infra、编译原理和本地推理引擎放到一张图里。读者需要先知道全书在建什么系统、每层交付什么、C++20 项目应该如何分模块，后面才不会把 manifest、IR、kernel、runtime 混在一起。

\topic{第二章：张量、算子、量化与计算账本。} 从张量布局、地址公式、matvec、GEMM、packing、SIMD、量化合同和 7B 账本开始。这章的目标是让读者能看懂一个 kernel 为什么快或慢：不是看函数名，而是看地址流、累加器、对齐、尾块、scale 和字节账本。

\topic{第三章：Transformer 推理、KV cache 与采样。} 把 decoder-only Transformer、reference decoder、prefill/decode、KV cache、logits 和 sampler 接成一条推理链路。这里开始把单算子性能放回整机系统。

\topic{第四章：AI 编译器与推理运行时。} 讲模型导入、manifest、schema resolution、shape verifier、dtype verifier、memory planner、graph/tensor/loop IR、融合、backend lowering 和 executable plan。读者不需要先懂传统编译器，本册会用推理引擎对象逐步解释每个概念。

\topic{第五章：C++20 CPU/GPU 后端优化。} CPU 后端要细化到 C++20 模块划分、RAII 资源管理、\code{std::span} 张量视图、aligned allocator、packing cache、SIMD micro-kernel、线程池、NUMA 放置、perf 计数器和 benchmark harness；GPU 后端要细化到 device buffer、stream、kernel launch、host-device staging、workspace 和 fallback；所有优化都必须配 oracle、profiler 和可复现实验。

\topic{第六章：工程验收、路线图与作品集证据。} 最后把测试、coverage、sanitizer、benchmark、perf 报告、框架对标、负例诊断和阶段路线图收束成作品集证据。没有这些证据，书稿再厚也只能证明学习过，不能证明能做生产级 AI Infra。

\subsection{C++20 项目应该怎样分层}

本册贯穿项目用 C++20 实现时，工程目录不能堆成一个巨大的 \code{engine.cpp}。一个保守、可维护的分层可以这样设计：

\begin{lstlisting}[numbers=none,caption={C++20 AI inference infra 的建议模块边界}]
include/ai/
  model/        manifest, tensor index, tokenizer contract
  graph/        graph IR, tensor descriptors, verifier facts
  compile/      passes, fusion, memory planner, lowering
  runtime/      session, KV cache, arena, scheduler, sampler
  backend/
    cpu/        packed weights, SIMD kernels, thread scheduler
    gpu/        device buffers, streams, kernel adapters
  test/         reference, oracle, profiler report schema

src/
  model/
  graph/
  compile/
  runtime/
  backend/cpu/
  backend/gpu/
  test/
\end{lstlisting}

这些目录不是形式主义。每个目录对应一组不变量：

\begin{itemize}
  \item \code{model/} 只把外部文件变成可信 manifest，不选择 kernel。
  \item \code{graph/} 表达语义和 shape/dtype 合同，不知道 AVX 寄存器有多少个。
  \item \code{compile/} 做 pass、memory plan 和 lowering，不直接持有某次用户请求的 KV cache。
  \item \code{runtime/} 绑定 session 状态并调度 executable plan，不重新解释模型文件。
  \item \code{backend/cpu/} 和 \code{backend/gpu/} 接收已经验证过的 descriptor，输出可测的 kernel 结果和 profiler 事件。
\end{itemize}

C++20 代码应该优先使用明确所有权和轻量视图：权重文件用 RAII 对象管理映射或加载缓冲，张量视图用 \code{std::span} 和 descriptor 表达，错误用明确结果类型返回，pass 用小对象或函数组合，热路径避免字符串查表、重复分配和不必要拷贝。后续写到实现时，类成员访问、常量命名、无递归遍历、模块分解和测试覆盖都要按同一套工程纪律执行。

\subsection{后端优化的总路线}

后端优化不能从“写 SIMD”开始。正确路线是先建立 reference，再建立可测 baseline，然后每次只改变一个机器事实。

\begin{enumerate}
  \item Reference：写最清晰的标量实现，定义正确输出和误差边界。
  \item Baseline kernel：用普通连续循环实现，记录吞吐、延迟、内存峰值和 profiler 事件。
  \item Layout 修正：确认 shape、stride、对齐、尾块和 cache line 行为。
  \item Packing：把权重从文件布局变成 kernel 友好布局，并记录 packed descriptor。
  \item SIMD micro-kernel：明确向量宽度、累加器数量、unroll、tail path 和精度。
  \item Thread scheduling：决定按行、列、token、head 还是 layer 切分任务，处理 barrier 和 false sharing。
  \item NUMA/cache/TLB：控制内存放置、预取、页行为、线程亲和和大页策略。
  \item Runtime integration：把 kernel 接进 prefill/decode plan，不让单算子优化破坏整体延迟。
  \item Oracle/profiler：每一步都比较误差和性能，不能只看最终 tokens/s。
\end{enumerate}

这条路线会贯穿后续内容。只要某个优化不能说明它改变了哪条地址流、哪个生命周期、哪个同步点或哪个后端 layout，它就还不是一个合格的工程优化。

\subsection{本节验收线}

读完本节，应该先建立全书地图：

\begin{itemize}
  \item 编译器的核心是理解语义、建立中间表示、检查不变量、在语义不变下优化机器执行。
  \item AI 编译器把模型资产、张量合同、量化元数据、推理图和后端能力编成 executable plan。
  \item 推理引擎不是单个 kernel，而是模型资产层、编译前端层、IR 优化层、后端层和运行时层的组合。
  \item C++20 项目要按 model、graph、compile、runtime、backend、test 分层，避免大文件和跨层猜测。
  \item 后端优化必须按 reference、baseline、layout、packing、SIMD、线程、NUMA、runtime、oracle/profiler 的顺序建立证据。
\end{itemize}

后面的章节都要回到这张地图。读者即使零基础，也可以按“这一章属于哪一层、建立哪个合同、优化哪条机器账本、怎样验证”这四个问题往下读。量化相关章节还要多问一个问题：低比特表示是否同时保持了数值质量、metadata 一致性和后端吞吐。
